% Death Spiral: Purity decay below P_crit = 2/7
% ----------------------------------------------------------------
% Mathematical content (Theorem, Eq. eq:death-spiral):
%   || Gamma(tau) - I/7 ||_F  <=  || Gamma(0) - I/7 ||_F * exp(-Delta*tau)
% Using || Gamma - I/7 ||_F^2 = P - 1/7 (derived in App. A, derivation 1):
%   ( P(tau) - 1/7 )  <=  ( P(0) - 1/7 ) * exp( -2*Delta * tau )
% The purity deficit therefore decays with rate 2*Delta, NOT Delta.
%
% For illustration we use Delta = 0.2, so the purity deficit decays
% with rate 2*Delta = 0.4. All three trajectories use this common
% Delta so that the figure has a consistent time-scale.
%
% Trajectories:
%   Green "attractor"  : P settles at 3/7 (Goldilocks upper bound, T-124)
%   Blue  "perturbed"  : P dips toward P_crit at tau = 2 then recovers
%   Red   "death spiral": P starts at 0.27 (strictly below P_crit) and
%                         decays to 1/7 following the theorem bound
% ----------------------------------------------------------------

\begin{tikzpicture}
\begin{axis}[
  width=11cm, height=7cm,
  xlabel={Internal time $\tau$},
  ylabel={Purity $P(\tau)$},
  xmin=0, xmax=10,
  ymin=0.11, ymax=0.52,
  grid=both,
  grid style={gray!20},
  title={\textbf{Death spiral below $P_{\mathrm{crit}}$}},
  title style={font=\large},
  % Legend placed at north-east with semi-transparent fill so that
  % the green viable trajectory (settling at y = 3/7 ~ 0.43) remains
  % visible through the legend box.
  legend pos=north east,
  legend style={font=\scriptsize, fill=white, fill opacity=0.72, text opacity=1, draw=gray!50},
  ytick={0.143,0.286,0.3,0.4,0.429,0.5},
  yticklabels={$\tfrac{1}{7}$,$\tfrac{2}{7}$,$0.3$,$0.4$,$\tfrac{3}{7}$,$0.5$},
  every axis label/.style={font=\small},
]

% === P_crit threshold line (red dashed) ===
\addplot[thick, red, dashed] coordinates {(0, 0.2857) (10, 0.2857)};
\addlegendentry{$P_{\mathrm{crit}} = 2/7$}

% === Thermal death line (gray dotted) ===
\addplot[thick, gray, dotted] coordinates {(0, 0.1429) (10, 0.1429)};
\addlegendentry{$I/7$: thermal death, $P = 1/7$}

% === Viable trajectory, attracts to Goldilocks upper bound 3/7 ===
% Damped oscillation around the T-124 attractor P = 3/7.
\addplot[thick, green!55!black, domain=0:10, samples=160]
  {0.4286 + 0.018*sin(deg(1.4*x))*exp(-0.18*x)};
\addlegendentry{Viable: attractor at $P = 3/7$}

% === Perturbation that recovers (stability radius, T-104) ===
% Gaussian dip centred at tau=2, minimum P(2)=0.29 stays strictly
% above P_crit = 2/7 so regeneration can re-establish the attractor.
\addplot[thick, blue!75!black, domain=0:10, samples=200]
  {0.42 - 0.13*exp(-0.7*(x-2)^2)};
\addlegendentry{Perturbation within stability radius}

% === Death spiral: exponential decay from P(0)=0.27 < P_crit to 1/7 ===
% P(tau) - 1/7 = ( 0.27 - 1/7 ) * exp(-2*Delta*tau), with Delta = 0.2
% (illustrative value for Delta(L_0); canonically Delta > 0 is the
% spectral gap of the linear Lindbladian from T-39a).
\addplot[thick, red!70!black, domain=0:10, samples=200]
  {0.1429 + (0.27 - 0.1429)*exp(-0.4*x)};
\addlegendentry{Sub-threshold: death spiral}

% === Initial-condition markers (small dots) ===
\addplot[only marks, mark=*, mark size=1.6pt, green!55!black]
  coordinates {(0, 0.4286)};
\addplot[only marks, mark=*, mark size=1.6pt, blue!75!black]
  coordinates {(0, 0.399)};
\addplot[only marks, mark=*, mark size=1.6pt, red!70!black]
  coordinates {(0, 0.27)};

% === Regime annotation near the attractor (far LEFT, 2 lines) ===
% Narrow 2-line label placed at x=1.6 so that the widest legend entry
% ("Perturbation within stability radius") at north east cannot reach
% this label horizontally regardless of exact font metrics.
\node[font=\scriptsize, fill=white, fill opacity=0.9, text opacity=1,
      inner sep=1.5pt, text width=1.9cm, align=center]
  at (axis cs:1.6, 0.475)
  {Regeneration\\$\geq$ dissipation};

% === P-form exponential bound near the red curve (LEFT of the legend) ===
% The annotation matches the plotted axis (purity), NOT the raw
% Frobenius distance. This corrects the coefficient: the deficit
% (P - 1/7) decays at rate 2*Delta, not Delta. We explicitly set
% draw=none and use text=... so no visible red rectangle is produced.
\node[font=\scriptsize, fill=white, fill opacity=0.92, text opacity=1,
      inner sep=1.5pt, draw=none, text=red!70!black, align=center]
  at (axis cs:4.5, 0.21)
  {$P(\tau) - \tfrac{1}{7} \leq
    (P_0 - \tfrac{1}{7})\, e^{-2\Delta(\mathcal{L}_0)\,\tau}$};

% === Arrow showing irreversible decay, now anchored ON the red curve ===
% At tau=1 the red curve is at 0.1429 + 0.1271*exp(-0.4) = 0.228
% At tau=3 the red curve is at 0.1429 + 0.1271*exp(-1.2) = 0.181
% Arrow goes along the curve between these points.
\draw[-{Stealth[length=5pt]}, thick, red!60!black]
  (axis cs:1.0, 0.228) .. controls (axis cs:2.0, 0.200)
  .. (axis cs:3.0, 0.181);
\node[font=\scriptsize, color=red!60!black]
  at (axis cs:2.35, 0.243) {irreversible};

\end{axis}
\end{tikzpicture}
